\newcommand{\NoMargin}{0.10717}
\newcommand{\YesMargin}{0.11}
\newcommand{\ConceptTrainCount}{8}
\newcommand{\ConceptTestCount}{9}
\newcommand{\ParaphraseDescription}{GPT-6.1-sol generated one paraphrase for each of 100 fixed test sources (25 per dataset/class, seed 15), using the earlier override-paraphrase toy example as inspiration. No detector scores were used to select variants. Automated review judged intent preserved for 100/100, but agreed with only 83/100 publisher labels under the narrower override definition. The disagreements concern original positive examples whose standalone intent is ambiguous, such as a training-date question. Preserve publisher labels; these data do not establish independently verified semantic ground truth.}
\newcommand{\ParaphrasePlanStatus}{Done}
\newcommand{\ParaphrasePlanDate}{Nov. 20}
\newcommand{\GenerationDisclosure}{GPT-6.1-sol generated the 100 paraphrases and supplied automated intent/concept-pair reviews through the author-provided local API. Review prompts and all generated texts/judgments are retained in the data notebook and CSVs.}
\newcommand{\ThresholdDescription}{Frozen threshold: \texttt{0.37308203458513201}; seed-15 LoRA: \texttt{0.50968770950539799}; zero-shot: \texttt{0.4724093079566955}. LoRA seeds 16/17 use \texttt{0.4998234929255746}, \texttt{0.49177553307711619}, respectively.}
\newcommand{\RandomReferenceDescription}{Cells with AP at/below its random reference or AUROC at/below 0.5: frozen: \code{deepset\_base64}, \code{safeguard\_base64}, \code{safeguard\_zero\_width}; lora: \code{deepset\_base64}, \code{safeguard\_base64}, \code{safeguard\_zero\_width}; zero\_shot: \code{deepset\_base64}, \code{deepset\_homoglyph}, \code{deepset\_zero\_width}, \code{safeguard\_homoglyph}, \code{safeguard\_zero\_width}. Obfuscation/tokenization shift and widespread benign false alarms undermine the trained probes; the generic English NLI readout also struggles on this mixed-language task. These are explanations to investigate, not established causes. Frozen and LoRA FPR reach 1.0 on both base64 sets and almost all zero-width cases; high injection recall there is not useful resistance.}
\newcommand{\DriftDescription}{LoRA mean relative drift ranges from 0.0402 to 0.0505 across datasets/seeds. Point-estimate verdicts follow the stated rule; intervals and seed variation must also be considered. A borderline yes whose interval crosses 0.05 does not establish a stable substantial effect. Positive drift establishes movement relative to the frozen reference, not improved detection.}
\newcommand{\ConceptDescription}{For the frozen mean direction, the held-out gap is 3.168 clean and 1.106 under homoglyphs. LoRA homoglyph gaps range 1.063--1.073; mean-direction changes range 0.0007--0.0022. The concept norm is 5.459 frozen versus 5.396--5.458 after LoRA. Neither a larger concept norm nor a stable direction alone demonstrates robustness.}
\newcommand{\FigureFinding}{LoRA changes deepset homoglyph AUROC from 0.7768 to 0.7738 (-0.0030), below the practical improvement threshold; the confidence intervals show substantial uncertainty.}
\newcommand{\ClaimDescription}{Under publisher labels, this one-epoch LoRA pilot produces no practical AUROC improvement in the 8 evaluated perturbation columns under the +0.11 rule, despite measurable representation drift; this does not prove equivalence or rule out stronger PEFT schedules.}
\newcommand{\ErrorSampleDescription}{The sample contains 10 false positives and 2 false negatives.}
\newcommand{\ErrorFinding}{The sampled false alarms include ordinary German and English requests, while the missed injections include explicit and implicit redirection. This points to investigating language balance and tokenization alongside override semantics, rather than treating a high clean AUROC as semantic validation.}
\newcommand{\BaselineReproduction}{Clean AUROC reproduces the midsemester results to four decimals; maximum probability difference is 0.000884, so thresholds are recalibrated on validation rather than assuming bitwise CPU/GPU agreement.}
\newcommand{\TruncationDescription}{At the shared 128-token encoder budget, changed-text truncation counts are base64 145/315; homoglyph 85/315; paraphrase 4/100; zero\_width 116/315. Truncation is therefore part of the measured stress condition, not an isolated semantic perturbation effect.}
\newcommand{\ComputeDescription}{The measured LoRA training mean is $B=0.175$ min and extraction mean $E=0.399$ min. For three seeds and three epochs, provisionally budget $9B$ each for adapters and prompt/prefix tuning; allow $72B$ for full fine-tuning (eight times the restricted pilot's epoch cost), $9E$ for extraction across the nine added runs, and 4 min for lexical/external checks and visualization. Total: $90B+9E+4=23.4$ T4 minutes. API-generation latency and manual review are additional wall time.}
